This section details the data construction pipeline for \benchmark, which is illustrated in Figure~\ref{fig:framework}.

\subsection{High-quality Data Crawling and Filtering Stage}
To ensure high-quality and large-scale data collection, we follow a multi-step data mining process.

\subsubsection{Representative Repository and Version Selection}
\label{sec:repo_selection}
To construct a robust benchmark, we follow the repository selection strategy established by SWE-Bench~\cite{swebench}, a widely recognized issue-resolving benchmark encompassing repositories from diverse application domains with active development communities. Specifically, we use 12 repositories from SWE-Bench as the sources for PR metadata collection. For each repository, we designate a single snapshot as the evaluation target: the latest stable version available in SWE-Bench (2023). Thus, our benchmark comprises \textbf{12 distinct repository versions} (one per repository) for documentation generation.

\subsubsection{Large-scale PR Metadata Collection}
We leverage the GitHub REST API~\cite{githubapi} to crawl all PRs from the selected repositories to ensure a reliable collection. For each PR, we systematically extract and structure essential metadata, such as its unique ID, title, description, and code changes. This initial collection process yields a massive corpus of 177.4k PRs, which are then stored in a structured format for subsequent filtering.

\subsubsection{Multi-dimensional Data Filtering}
Raw crawled data is often noisy. To guarantee the high quality of our benchmark, we apply a multi-step filtering process.

\textbf{\textit{Basic Regularity Filtering:}} This step aims to quickly filter noisy PRs from the initial dataset. The specific rules are as follows:
\begin{itemize}[leftmargin=*]
    \item \textbf{Status Check:} Retain PRs that have been merged, as these represent accepted contributions that meet repository standards.
    \item \textbf{Milestone Tag:} Retain PRs containing the ``milestone'' attribute, as these are typically associated with major repository goals.
    \item \textbf{Length Constraint:} Retain PRs with description length greater than 50 characters, ensuring each PR provides clear context.
    \item \textbf{Branch Filter:} Retain PRs merged into the main branch, as these are formal contributions central to the repository's evolution.
    \item \textbf{Review Presence:} Retain PRs with review comments, indicating that the changes have undergone human validation.
    \item \textbf{Bot Exclusion:} Exclude PRs generated by bots to focus on human-made changes.
    
\end{itemize}

\textbf{\textit{Functionality Relevance Filtering:}}
This step aims to retain PRs focused on functionality implementation. 

\begin{itemize}[leftmargin=*]

    \item \textbf{Functionality Label:} Retain PRs that contain functionality-related labels, such as ``new feature'' or ``new API'', to ensure the relevance to functional contributions.
    \item \textbf{Functionality Modified File:} Retain PRs based on the filenames of the modified files. Specifically, a PR is kept if it modifies at least one functional file (such as ``.py'' extension) located within a functional directory (excluding directories like ``\texttt{/test}'').
    \item \textbf{Functionality Persistence:} Functionalities introduced by previous PRs may be altered or removed in subsequent code updates. To ensure the persistence of QA tasks, we focus on PRs merged before the snapshot time of the selected repository version in Section~\ref{sec:repo_selection}. We then verify the persistence of their introduced functionalities by checking whether the added non-comment code lines still exist in the corresponding files of the selected version, accounting for potential file renames.


\end{itemize}
After this series of rigorous filtering steps, we obtain 4,170 high-quality PRs, forming the comprehensive basis of our benchmark.

\subsection{Repository-level Context Retrieving Stage}
\label{sec:stage2}

Current benchmarks often assess in a fragmented manner, failing to evaluate overall accuracy. To address this, this stage aggregates abundant repository-level context for each PR, forming a solid foundation for constructing high-quality QA tasks.

\subsubsection{Program Dependency Analysis}

Code changes in PRs often propagate their impact beyond the immediately modified snippets, potentially affecting the dependent modules. Since raw diffs are fragmentary, we first utilize Tree-Sitter~\cite{tree-sitter} to parse and extract complete definitions of modified code snippets (e.g., functions, methods). To capture repository-level context, we analyze dependencies by identifying both the callers and callees of these snippets. 

\subsubsection{Associated Issue Retrieval}
Associated issues reflect repository requirements or feature motivations driving PR changes. We apply keyword-based regular expressions (e.g., \texttt{``closes''}) to extract related issue numbers from PR titles and descriptions, and then crawl issue metadata via the GitHub REST API. For the ``Django'' repository, where issues are tracked on its official website, we implement a custom crawler to ensure comprehensive coverage. This process grounds QA tasks in the broader repository context.

\subsubsection{External Web Page Extraction}
External web links in PR descriptions provide additional information, such as official documentation or community discussions. Using the BeautifulSoup package~\cite{BeautifulSoup}, we parse these pages and extract relevant information, supporting a more comprehensive understanding of the PR.

\subsubsection{Commit History Tracking}
PR typically consists of a series of commits, each documenting incremental changes. For every involved commit, we collect detailed metadata including commit message, code changes, and associated review comments. This information offers a holistic view of functionality evolution, deepening the overall understanding of PRs.

\subsection{Functionality-driven QA Construction Stage}

Current benchmarks often use vague scoring criteria, leading to unreliable evaluation. To overcome this, we leverage the rich PR context to construct functionality-driven tasks, enabling objective evaluation by comparing the LLM's answers with references. The core philosophy is to simulate documentation-driven development: a developer formulates a functionality requirement and consults the documentation for answers. Thus, each QA consists of two parts: 

\noindent \textbf{$\blacktriangleleft$} \textbf{\textcolor{violet}{Question}}: A composite input (developer's requirement) containing a \textbf{Functionality Description} and a \textbf{Query} (e.g., ``Determine if the functionality is implemented in the current repository?'').

\noindent \textbf{$\blacktriangleright$} \textbf{\textcolor{orange}{Answer}}: The factual answer extracted from the PR's metadata and context. Crucially, the functionality description serves as the content of the question, and LLMs should use this description to reason over the given documentation and predict the answer.

\subsubsection{Intent-guided Functionality Description Generation}
To generate high-quality functionality descriptions, we leverage LLMs' advanced contextual understanding capabilities. Specifically, we populate a predefined prompt template with the PR metadata and rich context from Section~\ref{sec:stage2}. We further employ the Chain-of-Thought (CoT) strategy~\cite{chain}, guiding the LLM to consider from a global perspective, which includes motivation, implementation details, and impact scope. The generated descriptions are structured along three intent dimensions to include technical details:
\begin{itemize}[leftmargin=*]
    \item \textbf{WHAT:} Entities constituting or affected by the functionality.
    \item \textbf{WHY:} Purpose and motivation behind the functionality.
    \item \textbf{HOW:} Detailed approaches used for implementation.
\end{itemize}
Here, \textbf{WHAT} and \textbf{HOW} can be derived from code changes and program dependencies, while \textbf{WHY} can be informed by commit messages and associated issues. Other contextual information further enriches the generated descriptions. 
To prevent answer leakage, the LLM is instructed to avoid explicit mentions of file paths or repository versions. Overall, the generated description simulates the information developers seek in the documentation.

\subsubsection{Reliable QA Pair Generation}
Based on the generated functionality descriptions, we formulate three kinds of QA tasks.

\paragraph{\textbf{Functionality Detection}}
This task focuses on detecting the presence of the described functionality in the current repository.
\begin{description}[leftmargin=0pt, labelwidth=2em]
    \item[\textbf{$\blacktriangleleft$} \textbf{\textcolor{violet}{Question}}:] We formulate the following question by combining the functionality description with a query:
\begin{myfindbox}
    \small \textit{\textbf{[Functionality Description]} + Determine if the functionality is implemented in the current repository?}
\end{myfindbox}
    \item [\textbf{$\blacktriangleright$} \textbf{\textcolor{orange}{Answer}}:] We compare the PR's merged time with the snapshot time of the selected repository version. The answer is \textbf{True} if the attribute is earlier than or the same as the target version, indicating the functionality is present; otherwise, the answer is \textbf{False}.
\end{description}

\paragraph{\textbf{Functionality Localization}}
This task focuses on locating the files responsible for implementing the described functionality.
\begin{description}[leftmargin=0pt, labelwidth=2em]
    \item[\textbf{$\blacktriangleleft$} \textbf{\textcolor{violet}{Question}}:] We formulate the following question by combining the functionality description with a query:
\begin{myfindbox}
    \small \textit{\textbf{[Functionality Description]} + Identify the code file(s) responsible for implementing the functionality?}
\end{myfindbox}
    \item [\textbf{$\blacktriangleright$} \textbf{\textcolor{orange}{Answer}}:] We generate the reference answer by extracting files from the PR's code changes, retaining functional files (such as \texttt{``.py''} extension) with newly added lines that are not located in non-functional directories (such as \texttt{``/docs''} or \texttt{``/tests''}).
\end{description}

\paragraph{\textbf{Functionality Completion}}
This task focuses on filling in the masked technical details of the described functionality.
\begin{description}[leftmargin=0pt, labelwidth=2em]
    \item[\textbf{$\blacktriangleleft$} \textbf{\textcolor{violet}{Question}}:] We construct the question by replacing technical details from the \textbf{WHAT}, \textbf{WHY}, and \textbf{HOW} dimensions of the functionality description with ``\texttt{[MASK]}'', and appending the query:

\begin{myfindbox}
    \small \textit{\textbf{[Masked Functionality Description]} + Fill in the \texttt{[MASK]} placeholders with the correct details?}
\end{myfindbox}
    \item [\textbf{$\blacktriangleright$} \textbf{\textcolor{orange}{Answer}}:] The reference answer is the list of corresponding technical details that are extracted from the functionality description.
\end{description}

As shown in Table~\ref{tab:benchmark_stats}, each entry in \benchmark features a detailed functionality description with an average length of 771.45 characters. On average, solving the questions in the entry requires locating 2.01 files and completing 7.48 details\footnote{Due to space limitations, the detailed prompt template and data structure are provided in our repository.}.

\subsubsection{Manual Quality Validation}
To validate the quality of \benchmark, we conduct a human calibration process. We randomly sample 100 entries and have them reviewed by two expert annotators, each possessing over three years of Python expertise. For the first task, annotators assess whether the existence of the described functionality aligns with the answer by carefully analyzing the code repository. For the second task, they verify that the answer (a file list) accurately corresponds to the described functionality through inspection of the code repository. For the third task, annotators ensure that the details required to fill are precise and can be sourced from PR's metadata and context. Besides, annotators evaluate whether the functionality description is clear and reasonable. Across all aspects, the inter-annotator agreement consistently exceeds a Kappa coefficient of 90\%, demonstrating high agreement.

\begin{table}[t]
\centering
\small
\caption{Statistics of the \benchmark. \textbf{\# Func. Desc.} denotes the average character length of functionality descriptions. \textbf{\% Detect. Ratio} is the percentage of positive entries for the detection task. \textbf{\# Loc. File} and \textbf{\# Comp. Detail} represent the number of files to locate and details to complete, respectively.}
\vspace{-1em}
\label{tab:benchmark_stats}
\renewcommand{\arraystretch}{0.6}
\setlength{\tabcolsep}{1.2pt}
\begin{tabular}{cc c ccc ccc}
\toprule
\multirow{2}{*}{\textbf{\# Entry}} & \multirow{2}{*}{\textbf{\# Func. Desc.}} & \multirow{2}{*}{\textbf{\% Detect. Ratio}} & \multicolumn{3}{c}{\textbf{\# Loc. File}} & \multicolumn{3}{c}{\textbf{\# Comp. Detail}}\\
\cmidrule(lr){4-6} \cmidrule(lr){7-9}
& & & Min & Max & Avg. & Min & Max & Avg. \\
\midrule
4,170 & 771.45 & 50.65 & 1 & 62 & 2.01 & 3 & 23 & 7.48 \\
\bottomrule
\end{tabular}
\vspace{-1.0em}
\end{table}
